\paragraph{Profiles below a core upper bound.}
For a fixed core operation $C$, consider
\[
 \mathcal L_c(C)=\{\Phi_c(H):H\text{ is core and }H\subseteq C\}.
\]
This is a finite distributive lattice: bounded helpers are closed under
owned choice, and choice by $c$ or by another owner gives profile
disjunction or conjunction, respectively. Its implication preorder can
already contain an arbitrary finite preorder on singleton coalitions.

\begin{proposition}[Preorders in principal profile images]
\label{core:prop:principal-profile-preorders}
Fix $K\geq3$ and distinct owners $a,b,c$. For every finite preorder
$(X,\preceq)$ there is a core operation $C$ on $X\sqcup\{p\}$ such that
\[
 \bigl\{\{x\in X:\Phi_c(H)(\{x\})=1\}:H\text{ core},\ H\subseteq C\bigr\}
 =\operatorname{Up}(X,\preceq),
\]
where $\operatorname{Up}$ denotes the upward-closed subsets. Restriction
to these singleton bits is a surjective lattice homomorphism from
$\mathcal L_c(C)$ onto this upset lattice.
\end{proposition}
\begin{proof}
For each upset $S\subseteq X$, put
\[
 \varphi_S(A)=[p\in A\text{ or }A\cap S\ne\varnothing].
\]
Take the greatest completion $E_S$ with $c$-profile $\varphi_S$,
$b$-profile $1-\varphi_S(V\setminus -)$, and unanimity at other
owners, where $V=X\sqcup\{p\}$. These profiles are admissible:
a positive $c$-coalition meets $S\cup\{p\}$, and a positive
$b$-coalition contains that set. Set $C=\sum_SE_S$.
This is a nonempty finite core union, and $E_S\subseteq C$ realizes
exactly $S$ on the singleton bits indexed by $X$.

If $u\preceq v$ and $u\ne v$, use the input with $u=\{a\}$,
$v=\{c\}$, and every other input, including $p$, equal to $\{b\}$.
Here $E_S$ returns $c$ if $v\in S$, and otherwise returns $b$:
upward closure then also gives $u\notin S$. The empty and full upsets
show that $C$ is exactly $\{b,c\}$ at this input.
The pair test \eqref{core:eq:profile-pair-safe} for
$\Phi_c(\{u\})=1$, $\Phi_c(\{v\})=0$ fails under this bound.
Indeed the difference input is $\{a\}$, the $c$-atomic set is
$\{v\}$, and the $b$-atomic set misses $u$.
Every bounded helper therefore respects $u\preceq v$. This proves the
equality, and evaluation of profile bits preserves both lattice operations.
The padding variable also handles $S=\varnothing$ and $X=\varnothing$.
\end{proof}

The number of simultaneous value bounds needed for one such implication
can be exponential, even with trivial singleton profiles and pure closures.

\begin{theorem}[Exact exponential threshold for a principal profile implication]
\label{core:thm:principal-profile-locality}
For every fixed $K\geq3$ and $d\geq2$, there are a core operation $C$
on $2d$ inputs, an owner $c$, and coalitions $A,B$ such that:
\begin{enumerate}
\item Every core $H\subseteq C$ satisfies
  $\Phi_c(H)(A)\leq\Phi_c(H)(B)$.
\item For any sample of at most $2^d-2$ valuations, there is a core $H$
  respecting the $C$-bound at every sampled valuation and satisfying
  $\Phi_c(H)(A)=1$, $\Phi_c(H)(B)=0$.
\item A sample of $2^d-1$ explicit atomic valuations excludes every
  such separator.
\end{enumerate}
The baseline $C$ has unanimity singleton profiles and identity pure-owner
closures.
\end{theorem}
\begin{proof}
Use the paired transversals in \cref{core:thm:exponential-locality},
with $m=2^d$, cyclic indices $0,\ldots,m-1$, and incidence
$L_j\subseteq U_i$ exactly when $j=i$ or $j=i+1$.
Recall $L(w)=\{j:w_j=\{c\}\}$, $U(w)=\{j:w_j\ne\{b\}\}$,
and $\mathcal C_i=\{w:(L(w),U(w))=(L_i,U_i)\}$.
Let $C=F_0$: it returns $\{b,c\}$ on the exceptional regions
$\mathcal C_i$ for $1\leq i<m$, and input union elsewhere.
The deletion-patch construction makes $C$ core. Put
\[
 A=U_{m-1},\qquad B=L_1.
\]

Its exact own-$c$ profile image is
\begin{equation}\label{core:eq:patch-profile-image}
 \mathcal L_c(C)=
 \{\varphi:\varphi\text{ normalized monotone},\
       \varphi(U_i)=\varphi(L_i)\ (1\leq i<m)\}.
\end{equation}
For necessity, take the atomic input $q_i$ with $c$ on $L_i$, $a$ on
the unique point of $U_i\setminus L_i$, and $b$ elsewhere.
Its bound is $\{b,c\}$. The pair test
\eqref{core:eq:profile-pair-safe} excludes a positive $c$-bit on $U_i$
and a negative one on $L_i$. Monotonicity gives equality.
For sufficiency use the completion $E_\varphi$ with $c$-profile
$\varphi$, dual $b$-profile, and unanimity elsewhere, as in
\cref{core:lem:crown-patch}. On a retained exceptional region it is
$\{c\}$ if the common endpoint bit is one, and $\{b\}$ if it is zero.
Elsewhere provenance bounds it by the input union. This proves
\eqref{core:eq:patch-profile-image} on arbitrary menus and for all extra
owners.

Form the directed graph of Boolean inclusions with reversed covers
$U_i\to L_i$ for $1\leq i<m$. There is a path
\[
 U_{m-1}\longrightarrow L_{m-1}\longrightarrow U_{m-2}
 \longrightarrow\cdots\longrightarrow U_1\longrightarrow L_1.
\]
Every path from $A$ to $B$ uses all these $m-1$ reversed covers.
A path that ends at rank $d$ cannot visit ranks above $d+1$, since
only the selected reversed covers decrease rank. At ranks $d,d+1$,
the incidence law permits only the displayed backward progress;
unselected vertices cannot provide another downward step.
Thus (1) follows from \eqref{core:eq:patch-profile-image}.

For each $1\leq i<m$, delete reversal $i$ as well. The reachability
upset of $A$ is then a normalized monotone $\varphi_i$: it contains
$V$ and cannot contain $\varnothing$, since every reversed cover ends
at rank $d$. It has
$\varphi_i(A)=1$, $\varphi_i(B)=0$, constant on every remaining
reversed interval. Its completion $H_i=E_{\varphi_i}$ respects the
$C$-bound everywhere outside $\mathcal C_i$. These regions are disjoint,
so every valuation can exclude at most one of the $m-1$ helpers.
Any $m-2$ valuations leave a common helper, proving (2).
The $m-1$ atomic inputs $q_i$ enforce all retained endpoint equalities,
and therefore exclude every separator, proving (3).

Finally, exceptional outputs have two winners, so singleton outputs
occur precisely on constant atomic tuples. An exception removes only
winners outside $\{b,c\}$ and requires both $b$- and $c$-atomic inputs.
No pure-owner source probe can therefore lose its supplied winner.
Thus all pure-owner closures are identities.
\end{proof}
